\section{Pasch switches and the selected quotient}

A Pasch switch changes four triples on six points while preserving pair
incidences. Its effect on the integral column lattice is controlled by the
common saturation of the old and new four-block lattices.

\begin{theorem}[Pasch transport]\label{thm:pt}
Let $D,D'$ differ by a single Pasch (quadrilateral) switch. The two \emph{local
four-block lattices}---spanned by the four old, resp.\ four new, blocks that
change---agree at every odd prime, and localized at $2$ each has index two in
their common saturation. Consequently every odd-primary part of $\Delt$ is
preserved, the defect-order ratio $|\Delt(D')|/|\Delt(D)|$ lies in $\{1/2,1,2\}$
\textup(so one of $\Delt(D),\Delt(D')$ may already equal the other\textup), and
the minimum numbers of binary generators of $\Delt(D)$ and $\Delt(D')$ differ by
at most one.
\end{theorem}

\begin{proof}
Four old blocks on six points are replaced by four new blocks. After labeling
these points $0,\ldots,5$, take the ordered old triples
$012,034,135,245$ and new triples $013,024,125,345$. If $O,N$ are their
$6\times4$ incidence matrices, then
\[
N=\tfrac12 O H,\qquad
H=\begin{pmatrix}1&1&1&-1\\1&1&-1&1\\1&-1&1&1\\-1&1&1&1\end{pmatrix},
\qquad H^2=4I_4.
\]
Both matrices have Smith factors $1,1,1,2$, by elementary row and column
reduction. Hence their column lattices $M_O,M_N$ have the same rational span,
each has index two in its integral saturation $\widehat M$, and they agree
over every $\ZZ_p$ with $p$ odd.

Since $2\widehat M\subseteq M_O\subseteq H_3$,
also $\widehat M\subseteq H_3$.

Put $L_0$ for the lattice generated by all
unaffected blocks. The exchange groups $E_O=\widehat M/M_O$ and
$E_N=\widehat M/M_N$ are both $\ZZ/2$. The \emph{global} defects
are $\Delt(D)=\Hk/(L_0+M_O)$ and $\Delt(D')=\Hk/(L_0+M_N)$, formed only after
adjoining the unchanged $L_0$; both surject onto the common quotient
$\Hk/(L_0+\widehat M)$ with kernel a quotient of $E_O$, resp.\ $E_N$. Hence every
odd-primary part is preserved, the defect-order ratio lies in $\{1/2,1,2\}$, and
the minimum binary generator numbers differ by at most one. Indeed, for
$Q=H_3/(L_0+\widehat M)$ and either defect $F$, a kernel generated by at most
one element gives $d_2(Q)\le d_2(F)\le d_2(Q)+1$, where
$d_2(F)=\dim_{\FF_2}(F_2/2F_2)$.
\end{proof}
